\documentclass[11pt]{report}

\usepackage{fancyhdr} % Cabeceras de página
\usepackage{lastpage} % Módulo para añadir una referencia a la última página
\usepackage{titling} % No tengo claro para qué es esto
\usepackage[left=3cm,right=2.5cm,top=3cm,bottom=2cm]{geometry} % Márgenes
\usepackage[T1]{fontenc}
\usepackage[utf8x]{inputenc}
\usepackage{xspace}
\usepackage{graphicx}
\usepackage{tikz}
\usepackage{wrapfig}
\usepackage{hyperref}
\usepackage{amssymb}
\usepackage{multirow}
\usepackage[official]{eurosym}
\usepackage{enumitem}
\usepackage{pdfpages}
\usepackage{ifthen}
\usepackage{etoolbox} % Comandos guays.

\hypersetup{
  	hyperindex,
    colorlinks,
    allcolors=blue!60!black
}


\setcounter{secnumdepth}{2}
\renewcommand{\baselinestretch}{1.4}

\title{UAM Software Notification and Damage Management System \\ Fault Manager Lite \\ Risk Management and Planning Document}
\date{\today}
\author{{\Large Triforce} \\ \vspace{5pt} \textit{Iván Márquez Pardo, Víctor de Juan Sanz, Guillermo Julián Moreno}}

\fancyhf{}
\fancypagestyle{plain}{%
	\lhead{\raisebox{12pt}{\textsc{Risk Analysis} - \small Ref. TFC-UAM-01}}
	\chead{\centering \vspace{-15pt} \includegraphics[width =40 pt]{Logo.jpg}}
	\rhead{\raisebox{12pt}{\small Ver. 1.0 - \today \vspace{2pt}}}
	\cfoot{\thepage\ of \pageref{LastPage}}
	\rfoot{}
}

\newcounter{risks}[section]
\newcommand{\header}[1]{\\ \indent \textbf{#1}\hspace{10pt}}

\newcommand{\reqvertsep}{\vspace{-12pt}}

\newcommand{\riskcat}{\subparagraph{Category}}
\newcommand{\riskdesc}{\reqvertsep\subparagraph{Description}}
\newcommand{\riskprob}[1]{\reqvertsep\subparagraph{Probability} $#1 \%$}
\newcommand{\riskimpact}{\reqvertsep\subparagraph{Impact}}

\newenvironment{risk}[2][]{
	\ifthenelse{\equal{#1}{}}{}{\label{#1}}
	\refstepcounter{risks}
	\subsection{Risk \arabic{risks} - #2}
}{}

\makeindex

\begin{document}
\maketitle
\tableofcontents
\newpage
\chapter{Introduction}

\chapter{References}

\chapter{Project definition}

\chapter{Estimation, Planning and important requirements}

\chapter{Risk analysis}


% % % % % % % % % % Product engineering 

%
\begin{risk}{Change of the requirements}
\riskcat Requirements
\riskdesc As in every software project, requirements may be subject to change if the client identifies new needs or clarifies previously requested features.
\riskprob{30}.
\riskimpact Very low - 1. Given that the requirements have been already approved by the client in the requirement analysis document.
\end{risk}
%
\begin{risk}{Localization inside buildings}
\riskcat Design
\riskdesc The design of a system that can locate an user is a difficult task. In the design stage, we can find that this is actually an impossible tasks even when our preliminary analysis revealed that it cloud be done.
\riskprob{60}
\riskimpact High - 3.5. If the design can't be done, the requirements and design will need to be redone and there will be penalizations from the client.
\end{risk}
%
\begin{risk}{Changes in the interface with the mobile phones}
\riskcat Design
\riskdesc One of your target system are mobile phone operating systems. These are constantly changing and one of these changes may affect our project, either because of changed APIs or new or deprecated features.
\riskprob{70}
\riskimpact Low - 1. Even when it's probable that there's an update during the development of the project, it will probably consist of minor changes and won't break compatibility.
\end{risk}
%
\begin{risk}{Error-prone UAM's auth server interface}
\riskcat Design
\riskdesc We don't completely trust that the interface with the UAM authentication server is well defined and correct. A defensive programming approach may be necessary.
\riskprob{30}
\riskimpact Low - 2. If this risk materializes we will need to change our programming interfaces and extend the testing phase.
\end{risk}
%
\begin{risk}{Unknown required performance}
\riskcat Design
\riskdesc We haven't estimated the number of users that will access the system at once. This means that there's a possibility that our server becomes overwhelmed with requests and can't attend them all.
\riskprob{10} There's a low possibility that this kind of system becomes overloaded enough to require changes in the architecture or hardware.
\riskimpact Low - 2. If the system can't accommodate all the requests, a hardware upgrade will probably be enough.
\end{risk}


% % % % % Code and Unit Test

% Feasibility
%
\begin{risk}{Algorithm to assign task not defined}
\riskcat Design
\riskdesc We haven't defined the algorithm to automatically assign task to workers.
\riskprob{100} We will have to deal with it for sure.
\riskimpact{Medium - 3} If we don't prevent this risk, the programmers will have to spend time in defining this algorithm and the project may be delayed. 
\end{risk}
%
\begin{risk}{Algorithm to assign priority not defined}
\riskcat Design
\riskdesc We haven't defined the algorithm to automatically assign priority to reports.
\riskprob{100} We will have to deal with it for sure.
\riskimpact{Medium - 3} If we don't prevent this risk, the programmers will have to spend time in defining this algorithm and the project may be delayed. 
\end{risk}

% Testing
\begin{risk}{No test plan designed}
\riskcat Design 
\riskdesc We haven't designed the test plan.
\riskprob{100} We will have to deal with it for sure.
\riskimpact{High - 4} The programmers will have to spend time in defining the test plan and it may be not as extensive and exhaustive as a test plan done by designers. It may also affect to the project duration because testers will have to define the test plan on the go.
\end{risk}

% % % Integration and Test

% Environment
\begin{risk}{Difficult to test multi-user interaction}
\riskcat Design
\riskdesc  It will be difficult to develop realistic multi-user interaction for testing response requirements.
\riskprob{100} We will have to deal with it for sure.
\riskimpact{Very low - 1} We can apply statistics as we are used to do in SGNIP to estimate responses time in multi-user interaction scenarios.
\end{risk}

%Product

\begin{risk}{Acceptance criteria for requirements}
\riskcat  Design
\riskdesc We haven't agreed an acceptance criteria in a formal agreement for all requirements.
\riskprob{20} It is not very probable that UAM disagrees with us about the acceptance of requirements because of all the meetings we have had and the approve meanwhile.
\textcolor{red}{Check grammar}
% % % % % % % TODO:  Meanwhile¿¿?? I used gtraductor and i not very convinced.
\riskimpact{Low} Even if there is not a formal agreement, UAM has approved the requirements so it shouldn't cause us a big problem. 
\end{risk}

% % % % TODO: System i don't know what to say. (Questions 55-59)


% Maintainability
%\begin{risk}
%\riskcat Design
%\riskdesc Web developing can change within the years and that could difficult maintainability because maintenance people may have to change code to adapt the system to new specifications.
%\riskprob{20} When some new web standard is released it maintain compatibility with older systems.
%\riskimpact{}
%\end{risk}


% % % % % % % % % % % % % % % % % % Development Environment

%\begin{risk}
%\riskcat Development Environment
%\riskdesc We haven't defined requirements traceability that tracks requirements from source specification through test cases.
%\riskprob{70}  
%\riskimpact{}
%\end{risk}

% i jump from 117 to 142 because there are all questions about how we work in the company and i don't know anything and i don't want to invent either so, i go through.


% % % % % % % Program constraints. 
\begin{risk}{Subcontracted companies - Communication}
\riskcat Development Environment
\riskdesc There may be communication problems because we have subcontracted 2 company to help us with the development of the system.
\riskprob{30} We are used to deal with subcontracted companies so we know how to deal with communication issues. It is not very probable to happen.
\riskimpact{Medium - 3} When people working together can't communicate properly the development get delayed and looses quality.
\end{risk}

%
\begin{risk}{Subcontracted companies - Assigning task}
\riskcat Development Environment.
\riskdesc As we are working with 2 subcontracted companies, we will have to take care of assigning task and submodules to developers.
\riskprob{100} We will have to deal with it for sure.
\riskimpact{Low - 2} It is not a very big problem because the project leader has to assign task even if there are not subcontracted, so this just make his job a bit harder.
\end{risk}

%
\begin{risk}{Subcontracted companies - Integration test}
\riskcat Development Environment.
\riskdesc As we will not codify all submodules, integration test can get complicated.
\riskprob{50} If every worker sticks to the project plan, to the the analysis and to the design we should get too much work in integration test.
\riskimpact{Medium - 3} When integration test fails and we have to check someone else's code it get hard to get the integration done.
\end{risk}
%
\begin{risk}{Money and time limitation}
\riskcat Program constraints.
\riskdesc The timing has been limited in time and money. We may have to cancel the project.
\riskprob{75} Based on our experience, we are not positive that we can assume 10\% of the time and money 
\riskimpact{Very high - 5} The project may get canceled. 
\end{risk}


\begin{risk}{Inexperienced staff}
\riskcat Program constraints
\riskdesc The staff is inexperienced in this kind of applications and architecture and they may not work according to the plan.
\riskprob{70} According to our experience, when working with staff inexperienced it quite probable that the project get delayed.
\riskimpact{Medium - 3} The project may get delayed.
\end{risk}


\section{Risk Identification}
% -*- root: ../RiskAnalysis.tex -*-

This section contains all the generic and specific risks that have been identified taking into account the characteristics of the Fault Manager Lite project. Generic risks have been extracted from the answers to the questions of the Taxonomy-based Questionnaire contained in the Taxonomy-based Risk Identification document \cite{taxonomy93}, while the more specific risks of this project have been extracted after a deep analysis, evaluation and group discussion about the FML Project Plan \cite{plan15}.

Each risk has been given a numeric ID which will be used as a reference in the next sections of this document. Moreover, all the risks have been described and classified in categories, in order to better comprehend their scope.
Among these possible categories, there are:
\begin{itemize}
\item Related to requirements (R).
\item Related to the estimation (E).
\item Related to external sources (S).
\item Related to project conclusion (C).
\item Related to personnel management (P).
\item Related to supervision and monitoring (M).
\item Related to software development (D).
\end{itemize}

\begin{risk}{Frequent change of the requirements}
\label{riskReqChange}
\riskcat Requirements
\riskdesc As in every software project, requirements may be subject to change if the customer identifies new needs or clarifies previously requested features.
\end{risk}

\begin{risk}{Negative attitude of final users with respect to the system}
\label{riskAttitude}
\riskcat Project conclusion
\riskdesc A software application that is not used by its intended users after it has been released, must be considered as a failure.
\end{risk}

\begin{risk}{Unexpected necessary features}
\label{riskFeaturesUnexpected}
\riskcat Requirements
\riskdesc During the software development, some necessary functionalities that we have not taken into account might appear. If they are important or complex, they will affect the project planning in a significant way.
\end{risk}

\begin{risk}{Developed software does not meet customers' expectations}
\label{riskExpectations}
\riskcat Project conclusion
\riskdesc There has been some misunderstandings between the customer's intentions (expected requirements) and those that have been implemented by the development. Getting customer's validation requires re-work.
\end{risk}

\begin{risk}{Localization inside buildings}
\label{riskLocalization}
\riskcat Requirements
\riskdesc The design of a system that can locate an user is a difficult task. In the design stage, we can find that this is actually an impossible task even when our preliminary analysis revealed it was feasible.
\end{risk}

\begin{risk}{Error-prone UAM's authentication server interface}
\label{riskAuthServer}
\riskcat External sources
\riskdesc We don't completely trust that the interface with the UAM authentication server is well defined and correct. A defensive programming approach may be necessary.
\end{risk}

\begin{risk}{Changes in the interface with mobile phones}
\label{riskPhone}
\riskcat Software development
\riskdesc One of our target system are mobile phone operating systems. These are constantly changing and one of these changes may affect our project, either because of changed APIs or new or deprecated features.
\end{risk}

\begin{risk}{Lack of precedences}
\label{riskPrec}
\riskcat Software development
\riskdesc Our available technical staff does not have experience neither in this type of application nor in the underlying technical architecture. This fact may delay the development of the project.
\end{risk}

\begin{risk}{Difficulty on defining algorithms}
\label{riskAlgorithms}
\riskcat Requirements
\riskdesc We might encounter some difficulties while defining the two main algorithms of our software application: the automatic assignment of repair tasks to members of the maintenance service and the automatic assignment of priority to repair tasks.
\end{risk}

\begin{risk}{Problems with real time events/notifications}
\label{riskRealTime}
\riskcat Software development
\riskdesc We might encounter some difficulties while dealing with real time events and notifications. For example, when technicians update the status of repair tasks they have just done, managers reassigned tasks or users receive notifications/messages; these are real-time events that should be processed as soon as possible.
\end{risk}

\begin{risk}{Higher number of users than expected}
\label{riskUserLoad}
\riskcat Estimation
\riskdesc We might have underestimated the number of users that will access the system at once. This means that there's a possibility that our server becomes overwhelmed with requests and can't attend them all.
\end{risk}

\begin{risk}{First collaboration with subcontracted companies}
\label{riskCollaboration}
\riskcat Personnel management
\riskdesc Two companies are expected to be subcontracted: UAMSOFT Systems will participate in the development of the project, and SOFTCOM will take care of the software updates and product licenses that are necessary. We have never worked with any these companies, so we don't know their capabilities nor their working methods.
\end{risk}

\begin{risk}{System failures in design during integration tests}
\label{riskIntegrationTests}
\riskcat Project conclusion
\riskdesc As we are working with another two companies, misunderstandings may appear and cause problems during the integration tests. If these companies are assigned some modules to implement, once they complete them, it is likely that their integration becomes a problem that we didn't noticed during the design phase.
\end{risk}

\begin{risk}{Inefficient task assignment and coordination with subcontracted companies}
\label{riskAssignment}
\riskcat Supervision and Monitoring
\riskdesc If we don't receive some feedback about the development of the tasks assigned to the subcontracted companies, our project management will feel the effects of this lack of coordination, resulting on a bad scheduling and delays.
\end{risk}

\begin{risk}{Poor management and planning decisions}
\label{riskManagement}
\riskcat Supervision and Monitoring
\riskdesc Lack of communication and information flow causes a bad organization and coordination between members of the development team. Bad planning causes losses of time (workers are not assigned new tasks to perform, so they are waiting for them instead of working), and in last term, delays in the whole project.
\end{risk}

\begin{risk}{Delays in the calendar}
\label{riskDelays}
\riskcat Supervision and Monitoring
\riskdesc We might suffer from delays in the calendar with respect to the scheduling we had considered. These problems should be analyzed as soon as possible, in order to avoid that they propagate to other phases of the project.
\end{risk}

\begin{risk}{Excess of budget expenses}
\label{riskBudget}
\riskcat Supervision and Monitoring
\riskdesc We might suffer from excess of budget expenses with respect to the estimations we had considered. These problems should be analyzed as soon as possible, in order to avoid that they grow and multiply in the next phases of the project.
\end{risk}

\begin{risk}{Difficulties in maintaining the team united and motivated}
\label{riskMotivation}
\riskcat Project conclusion
\riskdesc In the last stages of the development, some quarrels or arguments may arise between some members of the team, while others might get demotivated. These issues may affect the rhythm of the project development.
\end{risk}

\begin{risk}{Poor quality of the product}
\label{riskQuality}
\riskcat Project conclusion
\riskdesc Due to time and budget restrictions, the quality of the final product might be affected, which means that a re-work is needed. This problem also applies to the deliveries planed on the set of milestones.
\end{risk}

\begin{risk}{High rotation of personnel}
\label{riskPersonnelRotation}
\riskcat Personnel management
\riskdesc The staff assigned to this project may suffer from rotations due to external causes that we can't control. We might have to deal with changes in the members of the development team, avoiding significant delays in the project duration.
\end{risk}

\begin{risk}{Underestimation of personnel needed}
\label{riskPersonnelUnderestimation}
\riskcat Estimation
\riskdesc During the development of the project, when revising the proportion of the project that we have done and the time we have spent on it, we might realize that we won't be able to finish it on time as a result of an underestimation of personnel needed.
\end{risk}

\begin{risk}{Non-specific assignment of responsibilities}
\label{riskResponsibilitesAssignment}
\riskcat Personnel management
\riskdesc If roles of the team members and authorities are not well defined in the early stages of the development, the responsibility for problems that appear later won't be easily assigned anybody. Also, if disagreements appear in the team, there won't be a single person to decide the best solution (leading to the implementation of several different solutions at a time).
\end{risk}


\section{Risk Estimation}
% -*- root: ../RiskAnalysis.tex -*-

In this section, we analyze each of the risks we have identified in the previous section. We have split the analysis into two subsections: in the first one, we estimate their probability of occurrence, while in the second we estimate the impact they have and consequences they may cause.

\subsection{Estimation of the probability of each risk occurrence}
The probability of occurrence of each risk has been valued from 0 to 1, attending to this scale:
\begin{itemize}
\item Very Low Probability: $0,1-0,2 (10\%-20\%)$.
\item Low Probability: $0,3-0,4 (30\%-40\%)$.
\item Medium Probability: $0,5-0,6 (50\%-60\%)$.
\item High Probability: $0,7-0,8 (70\%-80\%)$.
\item Very High Probability: $0,9 (90\%)$.
\end{itemize}


\begin{risk}[riskReqChange]{Frequent change of the requirements}
\riskcat Requirements
\riskprob{40} Low Probability

The FML project plan has already been approved by the customer. In this document, we completely specify and detail the functional and non-functional requirements of the application. As we received the approval from the customer, we don't expect many changes of the requirements, although we can't completely discard the possibility yet as the customers may come up with new ideas or changes as they start testing the deliverables.
\end{risk}

\begin{risk}[riskAttitude]{Negative attitude of final users with respect to the system}
\riskcat Project conclusion
\riskprob{20} Very Low Probability

In the Triforce company, we develop applications taking into account the intended final users and their opinions and suggestions we get from them during the demonstration of the application, near the end of project development. It is unlikely that this risk happens.
\end{risk}

\begin{risk}[riskFeaturesUnexpected]{Unexpected necessary features}
\riskcat Requirements
\riskprob{30} Low Probability

The FML project plan has already been approved by the customer. Despite this and the fact that we have carefully revised that document too, we are conscious that during the development we might encounter unexpected new features that are necessary and we had not realized they existed.
\end{risk}

\begin{risk}[riskExpectations]{Developed software does not meet customers' expectations}
\riskcat Project conclusion
\riskprob{10} Very Low Probability

In the Triforce company, we try to stay in contact with our customers and we also ask for their approval on interface mock-ups and prototypes. To date, we have never received a complain from a customer, but we don't discard this risk.
\end{risk}

\begin{risk}[riskLocalization]{Localization inside buildings}
\riskcat Requirements
\riskprob{70} High Probability

We have no first-hand experience with the GPS technology and its limitations, so it is likely that we expected too much from it when saying that it could give us the location of the users inside of the campus buildings.
\end{risk}

\begin{risk}[riskAuthServer]{Error-prone UAM's authentication server interface}
\riskcat External sources
\riskprob{50} Medium Probability

We don't really know the quality and security that the UAM's Authentication Server offers to those applications that connect to it. With no information about this issue, we assign it medium probability.
\end{risk}

\begin{risk}[riskPhone]{Changes in the interface with mobile phones}
\riskcat Software development
\riskprob{70} High Probability

Mobile phones are constantly updating and changing their specifications and interfaces, so it is quite likely that they change during the development of the project (and later, during the maintenance).
\end{risk}

\begin{risk}[riskPrec]{Lack of precedences}
\riskcat Software development
\riskprob{90} Very High Probability

It is highly likely that we will have to deal with problems either with the type of this application or with the underlying technical architecture, as our technical staff does not have any experience in these fields.
\end{risk}

\begin{risk}[riskAlgorithms]{Difficulty on defining algorithms}
\riskcat Requirements
\riskprob{30} Low Probability

Our technical staff has no experience with this particular type of application, but it should be noted that they have performed especially well on previous projects whose algorithms were more complex than the ones of this project. This risk is not likely to happen.
\end{risk}

\begin{risk}[riskRealTime]{Problems with real time events/notifications}
\riskcat Software development
\riskprob{20} Very Low Probability

Real time events or notifications are a quite delicate issue to deal with, but the amount of users connected at the same time to the application relaxes this condition of real time, lowering the probability of having problems with their performance.
\end{risk}

\begin{risk}[riskUserLoad]{Higher number of users than expected}
\riskcat Estimation
\riskprob{30} Low Probability

We are optimistic on the number of users the application will have. However, we don't expect that the number of people in the campus connected to our application at the same time, surpasses the estimation of 1,000 people. There's a low possibility that this kind of system becomes overloaded enough to require changes in the architecture or hardware.
\end{risk}

\begin{risk}[riskCollaboration]{First collaboration with subcontracted companies}
\riskcat Personnel management
\riskprob{80} High Probability

It is the first time that we work together with other two subcontracted companies, so at first it is highly likely that we can have problems defining common working procedures.
\end{risk}

\begin{risk}[riskIntegrationTests]{System failures in design during integration tests}
\riskcat Project conclusion
\riskprob{60} Medium Probability

Every subcontracted company will implement their own modules of the application and perform unit testing on them. As it is the first time we collaborate, we don't know their abilities, and the integration tests might fail as a result of a bad design or bad unit testing.
\end{risk}

\begin{risk}[riskAssignment]{Inefficient task assignment and coordination with subcontracted companies}
\riskcat Supervision and Monitoring
\riskprob{70} High Probability

The division of work into tasks that are the least related possible to each other might be quite complicated. Achieving it and dividing those tasks between us and the subcontracted companies will be useless without the delivery of feedback in order to supervise the advances of the project.
\end{risk}

\begin{risk}[riskManagement]{Poor management and planning decisions}
\riskcat Supervision and Monitoring
\riskprob{40} Low Probability

Our technical staff has previously worked together in other projects of Triforce, so they already know each other and we think that the information should easily flow. Project planning and task assignment would need more control as the duration of the project is tight. If personnel rotations happens, this risk will increase its probability.
\end{risk}

\begin{risk}[riskDelays]{Delays in the calendar}
\riskcat Supervision and Monitoring
\riskprob{80} High Probability

Coordinating efforts between three different companies is a difficult task that presumably will cause delays wit respect to the initial schedule. Also, it is the first time we work together so the probability of ocurrence of this risk rises.
\end{risk}

\begin{risk}[riskBudget]{Excess of budget expenses}
\riskcat Supervision and Monitoring
\riskprob{50} Medium Probability

As it is highly probable that the project may get delayed in greater or lesser extent, the workers will have to be working more hours and we have a material cost of 1,050\euro/month, which means we will spent more budget than the expected.
\end{risk}

\begin{risk}[riskMotivation]{Difficulties in maintaining the team united and motivated}
\riskcat Project conclusion
\riskprob{20} Very Low Probability

Our technical staff has previously worked together in other projects of Triforce: the information flow between members was acceptable and no big arguments were reported. However, we expect that the team might be under pressure (specially in late development stages), so we have to take this into account.
\end{risk}

\begin{risk}[riskQuality]{Poor quality of the product}
\riskcat Project conclusion
\riskprob{60} Medium Probability

We have a tight budget for the project, the duration of the project has been tightly limited and we don't know the abilities of the subcontracted companies. Taking into account these factors, it is likely that the quality of our final product gets affected.
\end{risk}

\begin{risk}[riskPersonnelRotation]{High rotation of personnel}
\riskcat Personnel management
\riskprob{30} Low Probability

High rotations of personnel are not that frequent in our company. Moreover, the duration of the project has been estimated in more or less 4 months and we don't have a large staff; therefore, the probabilities for a worker to leave the project are not specially high.
\end{risk}

\begin{risk}[riskPersonnelUnderestimation]{Underestimation of personnel needed}
\riskcat Estimation
\riskprob{10} Very Low Probability

We can only work with the technical staff that Triforce company has assigned to this project and we found it suitable when doing the project estimations. Also, we count with the help of two subcontracted companies, so it is unlikely that we lack of human resources.
\end{risk}

\begin{risk}[riskResponsibilitesAssignment]{Non-specific assignment of responsibilities}
\riskcat Personnel management
\riskprob{20} Very Low Probability

At the beginning of the development, roles and responsibilities will be clear, but as the project advances or personnel rotations happen, these may get diffused or even not really applied. The probability of ocurrence of this risk in a software project of the Triforce company is quite low, as roles and their responsibilities are defined from the very beginning of our projects; in case of personnel rotations, newcomers are assigned their roles as soon as they enter the project.
\end{risk}

\subsection{Estimation of the consequences of each risk}

The impact of each risk has been valued from 0 to 1, attending to this scale:
\begin{itemize}
\item Very Low Impact: 0,1.
\item Low Impact: 0,2.
\item Medium Impact: 0,3.
\item High Impact: 0,4.
\item Very High Impact: 0,8.
\end{itemize}

\begin{risk}[riskReqChange]{Frequent change of the requirements}
\riskcat Requirements
\riskimpact{0.3} Medium

If the requirements of the project are continuously changing due to numerous or complex requests from the customer, they will need changes in the schedule and they will possibly delay the whole project, specially if requests are received in the last development stages.
\end{risk}

\begin{risk}[riskAttitude]{Negative attitude of final users with respect to the system}
\riskcat Project conclusion
\riskimpact{0.3} Medium

If we don't make our application enough user-friendly and appealing to the users that they start and continue using it, it will have to be considered a failure. The reputation of our company might also be affected.
\end{risk}

\begin{risk}[riskFeaturesUnexpected]{Unexpected necessary features}
\riskcat Requirements
\riskimpact{0.4} High

If we realize that we have forgotten to develop some unexpected essential features, they might be a huge trouble to deal with, as if they are important they may require design changes and the later they are detected, the worse consequences they will have (mainly delays in the whole project).
\end{risk}

\begin{risk}[riskExpectations]{Developed software does not meet customers' expectations}
\riskcat Project conclusion
\riskimpact{0.4} High

The customer of the FML project in particular will receive a deliverable at the end of each development increment, being able to ask for interface or functionality changes before they grow and become more difficult to fix them.
\end{risk}

\begin{risk}[riskLocalization]{Localization inside buildings}
\riskcat Requirements
\riskimpact{0.2} Low

TODO
\end{risk}

\begin{risk}[riskAuthServer]{Error-prone UAM's authentication server interface}
\riskcat External sources
\riskimpact{0.3} Medium

Further investigation on this issue should be done before knowing the real impact it has on out project. If this risk materializes we will need to change our programming interfaces and extend the testing phase. With no information about this issue, we assign it medium impact.
\end{risk}

\begin{risk}[riskPhone]{Changes in the interface with mobile phones}
\riskcat Software development
\riskimpact{0.1} Very Low

Even when it is likely that there's an update during the development of the project, it will probably consist of minor changes and won't break compatibility.
\end{risk}

\begin{risk}[riskPrec]{Lack of precedences}
\riskcat Software development
\riskimpact{0.8} Very High

The later we start working on the project at 100\%, at our maximum abilities and performance, the more likely is that we are delayed with respect to the whole schedule and not being able to finish the project on time.
\end{risk}

\begin{risk}[riskAlgorithms]{Difficulty on defining algorithms}
\riskcat Requirements
\riskimpact{0.2} Low

In spite of having much experience defining algorithms, if we had underestimated their complexity and they required more time than the expected, our designer could be working on them while the programmers leave their implementation to the last moment of the coding phase. In the worst case, the algorithms could cause delays in the project schedule.
\end{risk}

\begin{risk}[riskRealTime]{Problems with real time events/notifications}
\riskcat Software development
\riskimpact{0.1} Very Low

Real time events are more a desirable characteristic than a mandatory one, so even if we find that we are having some troubles achieving a good performance on them, the impact will be low (for example, notifying that a task repair has been completed would take a whole minute instead of seconds, which is not very critical).
\end{risk}

\begin{risk}[riskUserLoad]{Higher number of users than expected}
\riskcat Estimation
\riskimpact{0.1} Very Low

If the system becomes so successful that many people collaborates with the maintenance service and our server can't handle all the requests, a hardware update or an increment on the server parallelism will probably be enough.
\end{risk}

\begin{risk}[riskCollaboration]{First collaboration with subcontracted companies}
\riskcat Personnel management
\riskimpact{0.3} Medium

We know nothing about those two companies. We don't know how they work or the capabilities they have. It is often difficult to coordinate workers in a single team, so coordinate efforts between three companies, it is even more difficult. Moreover, we have a tight delivery date, so we can't wait until perfectly knowing each other.
\end{risk}

\begin{risk}[riskIntegrationTests]{System failures in design during integration tests}
\riskcat Project conclusion
\riskimpact{0.3} Medium

Failures in integration tests will make us revise modules we have not implemented looking for mistakes instead of continuing with the schedule, delaying it in a significant way.
\end{risk}

\begin{risk}[riskAssignment]{Inefficient task assignment and coordination with subcontracted companies}
\riskcat Supervision and Monitoring
\riskimpact{0.8} Very High

Convenient assignment of tasks is crucial if we don't want to be working on the same thing twice or be unproductive. If we don't coordinate efforts with the other two companies, last minute delays may appear and cause a global delay on the project.
\end{risk}

\begin{risk}[riskManagement]{Poor management and planning decisions}
\riskcat Supervision and Monitoring
\riskimpact{0.4} High

Uncoordinated or useless efforts caused by lack of information feedback and/or supervisions result in last term on delays of the whole project. Personnel rotations might have great impact on the planning, as newcomers won't be assigned complex tasks until they know better the conditions of the project; revising the whole planning will be needed in this case.
\end{risk}

\begin{risk}[riskDelays]{Delays in the calendar}
\riskcat Supervision and Monitoring
\riskimpact{0.3} Medium

We should control the delays in order to limit them to early stages of the development, when we are still able to adapt future planning attending to this little delay at the beginning. If we can't achieve this, delays will grow and multiply and we won't finish the project on time.
\end{risk}

\begin{risk}[riskBudget]{Excess of budget expenses}
\riskcat Supervision and Monitoring
\riskimpact{0.4} High

If we are able to follow our initial time schedule and reach every milestone without any delays in the calendar, then our cost estimation should be correct. However, unexpected costs may appear (accidents, personnel rotation, time off sicks, broken hardware, etc), so they shouldn't be completely discarded.
\end{risk}

\begin{risk}[riskMotivation]{Difficulties in maintaining the team united and motivated}
\riskcat Project conclusion
\riskimpact{0.2} Low

If some members have big arguments, they will tend to lose their concentration on the project and their performance will be affected. The same happens to those members demotivated or depressed. This risk is only critical in late stages of the project development, when it is too late to change the planning.
\end{risk}

\begin{risk}[riskQuality]{Poor quality of the product}
\riskcat Project conclusion
\riskimpact{0.8} Very high

Delivering the customer a final product whose quality we know that it is quite poor, is quite frustrating and embarrassing. Also, the customer won't be happier either, never contracting our sevices again. Our reputation might get affected because of this risk.
\end{risk}

\begin{risk}[riskPersonnelRotation]{High rotation of personnel}
\riskcat Personnel management
\riskimpact{0.4} High

The same reasons that cause the low probability of ocurrence of this risk (tight duration of the project and small staff), are the causes of the high impact of this risk. The later a developer leaves the project, the more impact will have this risk on it, as we will need to find a substitute and the newcomer won't be able to do his/her best work from the beginnning, causing delays in the project that planning itself won't be able to avoid.
\end{risk}

TODO
\begin{risk}[riskPersonnelUnderestimation]{Underestimation of personnel needed}
\riskcat Estimation
\riskimpact{3} Medium

If we later detect that we can't reach the milestones on time even if our workers had been giving the best of them during the whole project; there is a lack of personnel (or the estimation was not accurate enough). We should then hire more people, which means more delays in the calendar (they have to catch up with the rest of the developers) and excesses on the budget (new unexpected salaries).
\end{risk}

\begin{risk}[riskResponsibilitesAssignment]{Non-specific assignment of responsibilities}
\riskcat Personnel management
\riskimpact{0.3} Medium

If two or more people are working on the same module, there should be a responsible that makes planning and design decisions. If no responsible is designated, contradictory solutions may be applied to the same problem causing delays in the schedule, as two or more people might have been working on the same issue.
\end{risk}


\chapter{Risk management}

\subsection{Risk 1 - Change of the requirements}
\subsection{Risk 2 - Localization inside buildings}
\subsection{Risk 3 - Changes in the interface with the mobile phones}
\subsection{Risk 4 - Error-prone UAM's auth server interface}
\subsection{Risk 5 - Unknown required performance}
\subsection{Risk 6 \& 7 - Algorithm to assign priority not defined}
We can prevent this risk by assigning personnel to defined the algorithm to automatically assign task and automatically assign priorities during the development phase so it get defined and designed when programmers get to develop that part in the first increment. 

The designer will have to work hard to get those algorithms defined in the first increment because is in the first increment when the task management will be developed.

\subsection{Risk 8 \& 13 - No test plan designed}
We have to define the test planning, taking into account that some parts of the system will be developed by an external company. 

We have to reestimate also the time and money spend in defining this test plan.

\subsection{Risk 9 - Difficult to test multi-user interaction}

\subsection{Risk 10 - Acceptance criteria for requirements}
\subsection{Risk 11 - Subcontracted companies - Communication}
\subsection{Risk 12 - Subcontracted companies - Assigning task}
\subsection{Risk 14 - Money and time limitation}
\subsection{Risk 15 - Inexperienced staff}




\chapter{Risk monitoring}


\chapter{Conclusions}

\appendix

\chapter{Riskology tool}

\end{document}
